\section{The question}
\label{s0:sec:question}

\subsection{The hypothesis in its author's words}

The starting point of this note is a reading of the rest-energy relation proposed by the author of the programme:
\begin{quote}
\emph{$E=mc^2$: consider $c^2$ as the symmetrical part of the equation and $m$ as what breaks this symmetry; $m$ is generated from
a $b$ particle, $b$ standing for the $b$-particle of Yang--Mills.}
\end{quote}
Three clarifications fix the intended reading. First, the energy is to be taken as \emph{one undivided entity}, and the mass is to
emerge through an \emph{internal differentiation} of that unity, before any external interaction is applied. Second, $c^2$ corresponds to
the variable $X$ of the former momentum--force chapter, whose successor is the interior dynamics of Part~IV of the paper
\cite{ES17} (Sections~15--17), where the shift operators are $\alpha_\pm=\partial_\tau\pm\kappa G$ with $\kappa=c/\ell$. Third, the energy
is six-dimensional. The $b$-field is the non-abelian gauge potential of Yang and Mills \cite{YangMills54}, whose quanta carry isotopic
spin; Yang and Mills left open whether those quanta are massive.

An earlier exploratory session settled several points that we take as given and do not repeat. The scalar $c^2$ is not by itself a
symmetry, and if $E$ and $c^2$ are invariant under a group then so is $m=E/c^2$, so a mass cannot \emph{be} the breaking of a symmetry
that $E$ and $c^2$ share. A nonzero invariant mass does not break Lorentz invariance. The free six-dimensional scalar has an energy
unbounded below (\cite{ES17}, Proposition~27.1), and positive four-dimensional mass cannot be obtained from momentum along the surplus
times, since $m_4^2=\mu_6^2-(p_u^2+p_v^2)/c^2$. Non-orthogonality of $\bE$ and $\bB$ generates no mass (\cite{ES17}, Section~14). The same
session was rightly criticised for jumping to Yang--Mills interactions, condensates and gap equations before the unity of the energy and
its differentiation had been built. This note does the work in the order the hypothesis asks for: the undivided energy first, its
differentiation second, the dynamics of the differentiation third, and the $b$-field last, where it belongs.

\subsection{The dictionary under test}

All constructions are made inside the arena of \cite{ES17}: $V=\R^{3,3}$ with $\eta=\operatorname{diag}(1,1,1,-1,-1,-1)$, the realisation
of $\Cl(3,3)$ on $\R^8$, the binding root $G$ ($Ge_i=e_{i+3}$, $Ge_{i+3}=-e_i$), which squares to $-\id$, is self-adjoint for $\eta$ and
acts as the Hodge star in the field-strength reading $X=\sum_i(B_ie_i+E_ie_{i+3})$, $z=\bB+i\bE$, $q(z)=z\cdot z$
(\cite{ES17}, Theorem~12.3). The working dictionary is
\begin{center}
\small
\begin{tabular}{@{}>{\raggedright\arraybackslash}p{3.4cm}>{\raggedright\arraybackslash}p{11.0cm}@{}}
\toprule
hypothesis & candidate object in the arena\\
\midrule
$E$, undivided, six-dimensional & the positive, duality-invariant energy $E_6=\tfrac12|z|^2=\tfrac12(|\bB|^2+|\bE|^2)$ of a frame vector; in
the closure of the force law, the Casimir $|p|^2$\\
$c^2$, the symmetric part & the conversion $C=c\,\id$ between the temporal and the spatial triple forced by Schur's lemma
(\cite{ES17}, Theorem~3.2), and the $\eta$-self-adjoint, factorising part $\kappa G$ of the shift, $(\kappa G)^2=-c^2/\ell^2$\\
$m$, the differentiation & the Lorentz-class invariant $E_0=\tfrac12|q|$, the rest-frame energy density of the field; for Dirac matter the mass
form $T=e_1e_2e_3$\\
$b$, the Yang--Mills field & the $\eta$-antisymmetric (transport) part of the shift, $\End(V)=\so{3,3}\oplus\{X=X^\dagger\}$
(\cite{ES17}, Theorem~17.3); at three constituents the binding-compatible algebra $\so4_\Delta\cong\su2\oplus\su2$ (\cite{ES17},
Theorem~17.21)\\
\bottomrule
\end{tabular}
\end{center}
Every entry is a hypothesis to be tested, and several of them turn out to need correction.

\subsection{How claims are classified}

We follow the deflationary principle of \cite{ES17} (Corollary~4.4): a statement about the arena carries information only if it could
have come out otherwise. Every claim below belongs to one of six classes, named in the text. A claim is \emph{fixed} when it is an identity
or theorem of the algebra that could have failed; \emph{constructed} when the algebra admits the object and nothing selects it;
\emph{obstructed} when it is impossible under stated hypotheses; \emph{open} when it is neither established nor refuted. To these four
classes of the paper we add two that the subject requires: \emph{inherited}, for facts of ordinary relativistic physics that hold in any
notation and are reproduced only because something else is stated in terms of them, and \emph{interpretive}, for readings and dictionary
entries, which are not theorems at all. The relation $E=mc^2$ itself is inherited. What can be new is only how the arena realises its
three ingredients, and which of the realisations are forced.

Every finite-dimensional statement was checked by exact or high-precision computation. Each section has a verification script in
\texttt{investigations/rest\_energy/verify/}, written independently of a second, adversarial check script by a different author; the
scripts print one \texttt{PASS} line per claim and exit with status zero only if all pass (\cref{s9:sec:methods}).
